\subsection{نتایج تجربی}
نتایج حاصل از ارزیابی تجربی بنچ‌مارک یکپارچه برای ماژول \lr{subspace\_projection} بر روی مدل‌های زبانی هدف در جدول \ref{tab:benchmark_subspace_projection} ارائه شده است:

\begin{table}[H]
\centering
\small
\caption{نتایج بنچ‌مارک تجربی ماژول تصویر زیرفضا \lr{subspace\_projection} روی مدل‌های پایه}
\label{tab:benchmark_subspace_projection}
\begin{tabular}{lccccc}
\toprule
\textbf{مدل پایه \lr{(Base Model)}} & \textbf{دقت \lr{(Accuracy)}} & \textbf{شاخص \lr{F1}} & \textbf{حافظه \lr{(Peak VRAM)}} & \textbf{زمان ارزیابی \lr{(Time)}} & \textbf{وضعیت} \\
\midrule
\lr{DistilBERT} & \lr{89.91\%} & \lr{0.8993} & \lr{338.8 MB} & \lr{18.5 s} & \lr{PASS} \\
\lr{RoBERTa} & \lr{84.17\%} & \lr{0.8535} & \lr{364.3 MB} & \lr{27.8 s} & \lr{PASS} \\
\lr{GPT-2} & \lr{59.52\%} & \lr{0.6596} & \lr{637.3 MB} & \lr{24.9 s} & \lr{PASS} \\
\lr{Qwen} & -- & -- & -- & -- & \lr{NOT\_RUN} \\
\bottomrule
\end{tabular}
\end{table}

یافته‌های آزمایشگاهی نشان می‌دهند که فشرده‌ساز \lr{subspace\_projection} دارای عملکردی متوازن و بسیار قوی در سراسر معماری‌های ارزیابی‌شده است؛ به طوری که در تمام مدل‌های \lr{DistilBERT}، \lr{RoBERTa} و \lr{GPT-2} به وضعیت قبولی قطعی \lr{PASS} دست یافته است. نکته حائز اهمیت، بهبود قابل‌توجه دقت در مدل \lr{RoBERTa} است که از دقت پایه \lr{75.11\%} به \lr{84.17\%} (+9.06\%) ارتقا یافته است؛ پدیده‌ای که ناشی از اثر تنظیم‌کنندگی (\lr{Regularization}) تصویر زیرفضای کم‌مرتبه و حذف نویز ابعاد فرعی وزن‌ها ارزیابی می‌شود. مدل \lr{Qwen} با وضعیت \lr{NOT\_RUN} مشخص شده است.
